%==============================================================================
% FICHAMENTO CRÍTICO — Diagnostic Performance of CNN-Based AI in Detecting
% Oral Squamous Cell Carcinoma: A Systematic Review and Meta-Analysis
% Conforme NBR 6023 (referências) e NBR 10520 (citações)
%==============================================================================
\section{Fichamento: \citeonline{shen2026cnn}}

% --- 1. IDENTIFICAÇÃO ---
\subsection{Identificação}
\begin{tabular}{ll}
\textbf{Título:} & Diagnostic performance of convolutional neural network-based
AI in detecting oral squamous cell carcinoma: a systematic review and
meta-analysis \\
\textbf{Autor(es):} & Shen, Mi; Jiang, Zhili; Feng, Yankun; Lin, Zhenzhen;
Lu, Cancan; Sun, Junli; Yao, Jun; Hu, Liang; Guo, Jincai \\
\textbf{Periódico:} & BMC Oral Health \\
\textbf{Ano:} & 2026 \\
\textbf{Volume:} & 26 \\
\textbf{Número:} & 1 \\
\textbf{Páginas:} & -- \\
\textbf{DOI:} & \url{https://doi.org/10.1186/s12903-025-07543-5} \\
\textbf{Qualis:} & A2 (Odontologia) \\
\textbf{Tipo:} & Revisão sistemática com meta-análise \\
\end{tabular}

% --- 2. OBJETIVO ---
\subsection{Objetivo}
Avaliar a acurácia diagnóstica da inteligência artificial baseada em redes
neurais convolucionais (CNN) na detecção do carcinoma espinocelular oral (OSCC),
por meio de revisão sistemática da literatura com meta-análise, sumarizando
sensibilidade, especificidade e razão de verossimilhança agrupadas.

% --- 3. METODOLOGIA ---
\subsection{Metodologia}
\textbf{Desenho do estudo:} Revisão sistemática com meta-análise conforme
diretrizes PRISMA-DTA. \\
\textbf{Amostra:} 14 estudos com 61.372 amostras no total, incluindo estudos
publicados até abril de 2025. \\
\textbf{Métricas principais:} Sensibilidade, especificidade, razão de
verossimilhança positiva (PLR), razão de verossimilhança negativa (NLR),
odds ratio diagnóstica (DOR), área sob a curva (AUC). \\
\textbf{Validação:} Análise com Meta-Disc 1.4 e Stata 18.0; avaliação de
heterogeneidade (I²); análise de subgrupos por desenho, sítio, método
estatístico, validação externa e tamanho amostral; nomograma de Fagan.

% --- 4. RESULTADOS PRINCIPAIS ---
\subsection{Resultados Principais}
\begin{itemize}
  \item A sensibilidade e especificidade agrupadas foram ambas de 0,94
        (IC 95\%: 0,89--0,98 e 0,92--0,97, respectivamente).
  \item A AUC foi de 0,98, com DOR de 261,58 (IC 95\%: 131,03--522,19).
  \item O PLR foi de 13,08 (IC 95\%: 9,21--18,60) e o NLR de 0,06
        (IC 95\%: 0,03--0,10).
  \item O nomograma de Fagan indicou que, com probabilidade pré-teste de 20\%,
        a probabilidade pós-teste aumenta para 81\%.
  \item Heterogeneidade significativa foi observada (I² > 75\%).
\end{itemize}

% --- 5. LIMITAÇÕES ---
\subsection{Limitações}
\begin{itemize}
  \item Alta heterogeneidade entre os estudos incluídos (I² > 75\%), não
        explicada totalmente pelas análises de subgrupo.
  \item Apenas um dos 14 estudos realizou validação externa, comprometendo
        a generalização dos achados.
  \item Os estudos são majoritariamente retrospectivos, com tamanhos amostrais
        limitados e variabilidade metodológica significativa.
\end{itemize}

% --- 6. CONTRIBUIÇÃO PARA O LIVRO ---
\subsection{Contribuição para o Livro}
Este artigo fundamenta o Capítulo 11, que aborda meta-análises de modelos CNN
para detecção de OSCC. Os dados agrupados (sensibilidade 94\%, especificidade
94\%, AUC 0,98) fornecem as estimativas mais robustas disponíveis até o momento
sobre o desempenho diagnóstico de CNNs no carcinoma oral. A heterogeneidade
identificada (I² > 75\%) e a carência de validação externa são discutidas no
livro como limitações críticas que precisam ser endereçadas pela metodologia
SDD/TDD proposta. O nomograma de Fagan (20\% $\rightarrow$ 81\%) é utilizado
para ilustrar o impacto clínico potencial da IA como ferramenta de triagem.

% --- 7. ANÁLISE CRÍTICA ---
\subsection{Análise Crítica}
\begin{itemize}
  \item \textbf{Validade interna:} Alta -- Protocolo PRISMA-DTA, busca em seis
        bases de dados, seleção por dois revisores independentes, análise com
        Meta-Disc e Stata.
  \item \textbf{Validade externa:} Média -- Incluiu estudos de múltiplos países,
        mas predominam dados asiáticos; apenas um estudo com validação externa.
  \item \textbf{Reprodutibilidade:} Alta -- Metodologia detalhada com estratégia
        de busca reproduzível e critérios de inclusão/exclusão explícitos.
  \item \textbf{Relevância clínica:} Alta -- Sensibilidade e especificidade de
        94\% indicam potencial para uso como ferramenta de triagem pré-biópsia.
  \item \textbf{Conflito de interesses:} Não -- Os autores declararam ausência de
        conflitos.
  \item \textbf{Viés identificado:} Viés de publicação -- a meta-análise incluiu
        apenas estudos publicados; a heterogeneidade elevada sugere possível
        viés de seleção de modelos com melhor desempenho.
\end{itemize}

% --- 8. CITAÇÕES RELEVANTES ---
\subsection{Citações Relevantes}
\begin{citacao}
``The pooled sensitivity and specificity of CNN based AI in detecting OSCC were
0.94 (95\% CI 0.89--0.98) and 0.94 (95\% CI 0.92--0.97). The area under the
curve being 0.98, respectively.'' (p. 10).
\end{citacao}

% --- 9. MAPEAMENTO SDD ---
\subsection{Mapeamento SDD}
\textbf{Problema:} Necessidade de estimativa robusta e agrupada do desempenho
diagnóstico de CNNs para detecção de OSCC, considerando a heterogeneidade entre
estudos e a carência de validação externa. \\
\textbf{Entrada:} 14 estudos primários com 61.372 amostras (imagens histopatológicas,
radiográficas e fotográficas), publicados entre 2018 e 2025. \\
\textbf{Saída:} Estimativas agrupadas de sensibilidade, especificidade, PLR,
NLR, DOR e AUC com intervalos de confiança de 95\%. \\
\textbf{Critérios:} Protocolo PRISMA-DTA; busca em $\geq$ 5 bases; qualidade
metodológica avaliada; meta-análise com modelo de efeitos aleatórios;
análise de heterogeneidade (I²). \\
\textbf{Confiança:} 0,80 -- Meta-análise bem conduzida, mas limitada pela alta
heterogeneidade e pela ausência de validação externa na maioria dos estudos
primários.

% --- 10. REFERÊNCIA ABNT ---
\subsection{Referência ABNT}
\begin{verbatim}
SHEN, Mi et al. Diagnostic performance of convolutional neural network-based
AI in detecting oral squamous cell carcinoma: a systematic review and
meta-analysis. BMC Oral Health, London, v. 26, n. 1, 2026. DOI:
10.1186/s12903-025-07543-5.
\end{verbatim}
